\section[再论 $\et(\Bun_G,\Lam)$]{再论 $\et(\Bun_G,\Lam)$（1 月 18 日）}


\srcpage{108}

设 $G/E$ 是约化群，$E$ 是剩余域为 $\Fq$、特征为 $p>0$ 的局部域。设 $\Lambda$ 是系数环，使得对与 $p$ 互素的 $n$ 有 $n\Lambda=0$。

回忆：$D_{\et}(\Bun_G,\Lambda)$ 有一个无限的半正交分解，分解为 $D_{\et}(\Bun^b_G,\Lambda)\cong D(G_b(E),\Lambda)$。

\subsection{紧对象}

回忆三角范畴 $C$ 中的对象 $A$ 是紧的（compact），是指 $\Hom_C(A,-)$ 与一切直和交换。

\begin{fact}
若 $C$ 是具备一切余极限的 $\infty$-范畴 $\cC$ 的同伦范畴，则 $A\in C$ 紧，当且仅当 $\Hom_C(A,-)$ 与一切余极限交换。
\end{fact}

\begin{remark}
在 abel 1-范畴中，紧性通常用与滤过余极限交换来定义。然而在三角范畴／稳定 $\infty$-范畴的情形下 $\Hom_C(A,-)$ 正合，故只要求与直和交换即可。
\end{remark}

\begin{fact}
若这样的 $C$ 由其紧对象 $C^\omega\subseteq C$ 在余极限下生成，则 $\mathrm{Ind}(C^\omega)\xrightarrow{\ \sim\ }C$ 是 $\infty$-范畴的等价。
\end{fact}

\begin{proposition}
$D(G_b(E),\Lambda)$ 是紧生成的，其紧生成元为 $\mathrm{c}\text{-}\mathrm{Ind}^{G_b(E)}_K\Lambda$，其中 $K\subseteq G_b(E)$ 是开 pro-$p$ 子群。
\end{proposition}

\begin{proof}
我们有
\[
\Hom_{G_b(E)}\bigl(\mathrm{c}\text{-}\mathrm{Ind}^{G_b(E)}_K\Lambda,A\bigr)=\Hom_K(\Lambda,A)=A^K.
\]
若 $K$ 是 pro-$p$ 的，则取 $K$-不变量（即 $K$-上同调）与一切直和交换。（若 $A$ 由复形 $A_\bullet$ 表示，则 $A^K$ 由逐层不变量 $A_\bullet^K$ 表示，该陈述可以显式验证。）故 $\mathrm{c}\text{-}\mathrm{Ind}^{G_b(E)}_K\Lambda$ 是紧的。为看出它们生成整个范畴，只需注意：若 $A$ 使一切 $K$ 都有 $A^K=0$，则 $A=0$。
\end{proof}

\begin{theorem}
范畴 $D_{\et}(\Bun_G,\Lambda)$ 是紧生成的；$A\in D_{\et}(\Bun_G,\Lambda)$ 紧，当且仅当对一切 $b\in B(G)$，记 $i^b\colon\Bun^b_G\hookrightarrow\Bun_G$，都有 $i^{b*}A\in D_{\et}(\Bun^b_G,\Lambda)\cong D(G_b(E),\Lambda)$ 紧，且对几乎所有 $b$ 消没。
\end{theorem}

\begin{proof}
先给出紧生成元。固定 $b\in B(G)$ 与开 pro-$p$ 子群 $K\subset G_b(E)$。目标是证明存在复形 $A^b_K\in D_{\et}(\Bun_G,\Lambda)$ 使得
\[
\Hom(A^b_K,B)=(i^{b*}B)^K.
\]
这样的对象由 Yoneda 引理刻画，且由定义显然（若存在）是紧的。这类对象（若存在）生成整个范畴。

为找出 $A^b_K$，我们使用上同调光滑的图卡
\[
\begin{tikzcd}
\wt\cM_b \arrow[d, "G_b(E)"'] \\
\cM_b \arrow[r, "\pi_b"'] & \Bun_G
\end{tikzcd}
\]
由此还得到 $f_K\colon[\wt\cM_b/K]\to\Bun_G$，它上同调光滑。

取 $A^b_K=Rf_{K!}f_K^!\Lambda$。为验证其有效，考虑
\[
R\Hom_{\Bun_G}(A^b_K,B)=R\Hom_{\Bun_G}(Rf_{K!}Rf_K^!\Lambda,B)=R\Hom_{[\wt\cM_b/K]}(Rf_K^!\Lambda,Rf_K^!B).
\]
\srcpage{109}

由 $f_K$ 对 $\Lambda$ 是 ULA 的（更确切地说，引理 18.25），可把后一表达式改写为
\[
R\Hom_{[\wt\cM_b/K]}(Rf_K^!\Lambda,f_K^*B\otimes Rf_K^!\Lambda).
\]
再利用 $f$ 的上同调光滑性（特别是相对对偶层 $Rf_K^!\Lambda$ 可逆），可把它改写为 $R\Gamma([\wt\cM_b/K],f_K^*B)$。最后，利用 $\wt\cM_b$ 的严格局部性把它等同于 $R\Gamma([\ast/K],f_K^*B|_{[\ast/K]})=(i^{b*}B)^K$。

我们还需检验所述紧对象的刻画。为此，对（拟紧）开子叠 $U\subset\Bun_G$ 作归纳。设 $A\in D_{\et}(\Bun_G,\Lambda)$ 紧。取 $b\in B(G)$ 使得 $i^b\colon\Bun^b_G\hookrightarrow U$ 是闭浸入（这之所以可能，只是因为我们截断到了拟紧的 $U$，否则会有无穷多个特殊化）。令 $j\colon V:=U\setminus\Bun^b_G\hookrightarrow U$。由归纳，已知 $D_{\et}(V,\Lambda)$ 的结论。

只需证明 $j^*$ 保持紧对象。的确，一旦知道这一点，若 $A$ 紧则 $j^*A$ 也紧，从而由正合三角形 $j_!j^*A\to A\to i^b_*i^{b*}A$ 知 $i^{b*}A$ 紧。所以 $j^*A$ 的紧性连同归纳给出：对一切满足 $\Bun^{b'}_G\subset U$ 的 $b'\ne b$，$i^{b'*}A$ 紧；而 $i^{b*}A$ 也紧。至于除有限多个之外这些茎必须消没，其证明与例 15.1 相同。反向类似（可由 $i^{b*}A$ 搭出 $A$；$j_!$ 保持紧对象是形式推理）。

现在证明 $j^*$ 保持紧性：只需在生成元上检验，而对「来自」$V$ 的生成元这一点是明显的。于是化归为检验 $j^*A^b_K\in D_{\et}(V,\Lambda)$ 紧。考虑交换图表（$U$、$V$ 都包含于 $\Bun_G$，图中略去包含映射）
\[
\begin{tikzcd}
{}[\wt\cM_b/K] \arrow[r, "f_K"] \arrow[d] & U \arrow[d, "j"] \\
{}[\wt\cM_b^\circ/K] \arrow[r, "f_K^\circ"'] & V
\end{tikzcd}
\]
其中 $\wt\cM_b^\circ=\wt\cM^\circ\setminus\ast$ 是有限 dim.tr.g. 的空间菱形。于是由 $A^b_K$ 的公式，$j^*A^b_K=Rf^\circ_{K!}Rf^{\circ!}_K\Lambda$。与上面类似的计算给出
\[
R\Hom(j^*A^b_K,B)\cong R\Gamma([\wt\cM_b^\circ/K],f^{\circ*}_K B)
\]
对一切 $B\in D_{\et}(V,\Lambda)$ 成立。可把它改写为 $R\Gamma(\wt\cM_b^\circ,B)^K$。但 $R\Gamma(\wt\cM_b^\circ,-)$ 与一切直和交换，因为 $\wt\cM_b^\circ$ 是有限 dim.tr.g. 的空间（因而 qcqs！）菱形\footnote{这一性质的一般证明我们尚未给出。它是有点棘手的点集拓扑。}。
\end{proof}

\begin{warning}
证明 $j^*A^b_K$ 紧时用到了 $j^*$ 的具体形式。并非每个开限制都保持紧性，因为 $\wt\cM_b^\circ$ 的一般开子集不会拟紧。
\end{warning}

\subsection{Bernstein–Zelevinsky 对偶}

我们将解释紧对象上的一个对偶。

\begin{proposition}
（Bernstein–Zelevinsky 对偶）对任意 $A\in D(G_b(E),\Lambda)^\omega$，存在唯一的 $A'\in D(G_b(E),\Lambda)^\omega$ 使得
\[
R\Hom(A',B)=(A\otimes B)_{hG_b(E)}.
\]
（导出同调函子 $(-)_{hG_b(E)}$ 是拉回 $D(\Lambda)\to D(G_b(E),\Lambda)$ 的左伴随。）

对 $A=\mathrm{c}\text{-}\mathrm{Ind}^{G_b(E)}_K\Lambda$，有 $A'=\mathrm{c}\text{-}\mathrm{Ind}^{G_b(E)}_K\Lambda$；一般地
\[
A'=R\Hom_{G_b(E)}(A,\cH(G_b(E))).
\]
双对偶映射 $A'':=(A')'\to A$ 是同构。
\end{proposition}

\srcpage{110}

\begin{proof}
由 Yoneda，$A'$ 若存在则唯一。为得存在性，只需取 $A=\mathrm{c}\text{-}\mathrm{Ind}^{G_b(E)}_K\Lambda$。此时有
\[
(A\otimes B)_{G_b(E)}=B_K\xrightarrow{\ \mathrm{Avg}\ }B^K=R\Hom\bigl(\mathrm{c}\text{-}\mathrm{Ind}^{G_b(E)}_K\Lambda,B\bigr).
\]
故所述 $A'$ 的公式在此情形有效。再利用这些 $A$ 生成整个范畴。
\end{proof}

\begin{theorem}
对任意 $A\in D_{\et}(\Bun_G,\Lambda)^\omega$，存在唯一的 $A'=\mathbb{D}_{\mathrm{BZ}}(A)\in D_{\et}(\Bun_G,\Lambda)$ 使得
\[
R\Hom(\mathbb{D}_{\mathrm{BZ}}(A),B)=\pi_\natural\bigl(A\otimes^L_\Lambda B\bigr),
\]
其中 $\pi_\natural\colon D_{\et}(\Bun_G,\Lambda)\to D(\Lambda)$ 是 $\pi^*$ 的左伴随（$\pi\colon\Bun_G\to\ast$）\footnote{如果你愿意，$\pi_\natural=R\pi_!(-\otimes^L R\pi^!\Lambda)$。}。

双对偶映射 $\mathbb{D}_{\mathrm{BZ}}(\mathbb{D}_{\mathrm{BZ}}(A))\to A$ 是同构。

若 $U\subset\Bun_G$ 是开子叠，则 $\mathbb{D}_{\mathrm{BZ}}$ 保持 $D_{\et}(U,\Lambda)^\omega\subseteq D_{\et}(\Bun_G,\Lambda)^\omega$。对 $U=\Bun^b_G$（$b$ 基本），它化归为 $D_{\et}(\Bun^b_G,\Lambda)^\omega=D(G_b(E),\Lambda)^\omega$ 上通常的 Bernstein–Zelevinsky 对偶。
\end{theorem}

\begin{remark}
更深层 $\Bun^b_G$（$b$ 非基本）上的层的 Bernstein–Zelevinsky 对偶会「展布」到开的层上。这与 Verdier 对偶形成对照：后者只向更特殊的层展布。
\end{remark}

\begin{proof}
我们在一类生成元上检验存在性。取 $i^b_!\bigl(\mathrm{c}\text{-}\mathrm{Ind}^{G_b(E)}_K\Lambda\bigr)$，其中 $i^b\colon\Bun^b_G\hookrightarrow\Bun_G$。我们断言 $\mathbb{D}_{\mathrm{BZ}}\bigl(i^b_!\mathrm{c}\text{-}\mathrm{Ind}^{G_b(E)}_K\Lambda\bigr)$ 就是 $A^b_K$。

为证此断言，我们检验
\[
R\Hom(A^b_K,B)=(i^{b*}B)^K\xrightarrow{\ \mathrm{Avg}\ }(i^{b*}B)_K=\pi_\natural\bigl(i^b_!\mathrm{c}\text{-}\mathrm{Ind}^{G_b(E)}_K\Lambda\otimes B\bigr).
\]
（原文在此式第二个等号处附注「up to shifts?」，即「差平移？」。）

为证双对偶，需要计算 $\mathbb{D}_{\mathrm{BZ}}(A^b_K)$ 是否等于 $i^b_![\mathrm{c}\text{-}\mathrm{Ind}^{G_b(E)}_K\Lambda]$。这在层 $\Bun^b_G$ 上容易看出，故只需在拉回到补集之后检验它。这就归结为如下陈述：左端在这样拉回之后消没。

令 $j\colon U\hookrightarrow\Bun_G$ 是由 $b$ 的真推广构成的开子叠。需要看出 $j^*\mathbb{D}_{\mathrm{BZ}}(A^b_K)=0$，等价地，对一切 $B\in D_{\et}(U,\Lambda)$ 有 $R\Hom(\mathbb{D}_{\mathrm{BZ}}(A),Rj_*B)=0$。按定义，这就是 $\pi_\natural(A^b_K\otimes Rj_*B)$。可用 $A^b_K$ 的公式计算它：
\[
\begin{aligned}
\pi_\natural(A^b_K\otimes Rj_*B)&\cong\pi_\natural f_{K!}(f_K^!\Lambda\otimes f_K^*Rj_*B)\\
&\cong(\pi\circ f_K)_!(f_K^*Rj_*B)\\
&\cong R\Gamma_c(\wt\cM_b/K,f_K^*Rj_*B).
\end{aligned}
\]
可把 $R\Gamma_c(\wt\cM_b/K,f_K^*Rj_*B)$ 看作 $[\wt\cM_b^\circ/K]$ 上朝 $[\wt\cM_b/K]$ 的边界紧支撑、且在 $[\ast/K]\hookrightarrow[\wt\cM_b/K]$ 附近不加支撑条件的上同调。这由「部分紧支撑上同调的消没」（第 19 讲的 Theorem 19.19）消失。
\end{proof}

\srcpage{111}

\subsection{Verdier 对偶}

Verdier 对偶是反变自函子 $A\mapsto R\Hom(A,R\pi^!\Lambda)$。它是 $D_{\et}(\Bun_G,\Lambda)$ 上的反变自函子。

$D_{\et}(\Bun^b_G)\cong D(G_b(E),\Lambda)$ 上的 Verdier 对偶就是光滑对偶（差平移与扭转）。

\begin{example}
$D(G_b(E),\Lambda)$ 的对偶复形是 Haar 测度模。
\end{example}

\begin{theorem}
对 $\Bun_G$ 的开子叠的任意开浸入 $j\colon U\hookrightarrow V$，以及任意 $A\in D_{\et}(U,\Lambda)$，有
\begin{enumerate}[label=(\arabic*)]
\item $Rj_*R\Hom(A,\mathbb{D}_U)\cong R\Hom(j_!A,\mathbb{D}_V)$；
\item $j_!R\Hom(A,\mathbb{D}_U)\cong R\Hom(Rj_*A,\mathbb{D}_V)$。
\end{enumerate}
\end{theorem}

\begin{proof}
(1) 由伴随性显然。

(2) 由归纳可设 $U=V\setminus\Bun^b_G$ 对某个 $b\in B(G)$ 成立，并令 $j\colon U\hookrightarrow V$。应用 $j^*$ 之后陈述显然。故只需证明：在作用 $R\Hom(A^b_K,-)$ 之后它仍是同构。由于 $R\Hom(A^b_K,B)=(i^{b*}B)^K$，左端消没。右端（差一个扭转）是
\[
R\Hom(A^b_K,R\Hom(Rj_*A,\Lambda))\cong R\Hom(A^b_K\otimes^L_\Lambda Rj_*A,\Lambda)\cong R\Hom(\pi_\natural(A^b_K\otimes^L_\Lambda Rj_*A),\Lambda).
\]
因此只需证明 $\pi_\natural(A^b_K\otimes^L_\Lambda Rj_*A)$ 消没。

利用 $\mathbb{D}_{\mathrm{BZ}}(A)=i^b_!(\mathrm{c}\text{-}\mathrm{Ind}^{G_b(E)}_K\Lambda)$，可把它改写为 $R\Hom(i^b_!(\mathrm{c}\text{-}\mathrm{Ind}^{G_b(E)}_K\Lambda),Rj_*A)$，而它消没，因为 $j^*i^b_!=0$。
\end{proof}

\begin{corollary}
$A\in D_{\et}(\Bun_G,\Lambda)$ 自反，即 $A\xrightarrow{\sim}\mathbb{D}(\mathbb{D}(A))$，当且仅当对一切 $b\in B(G)$，$i^{b*}A\in D_{\et}(\Bun^b_G,\Lambda)$ 自反，即对一切开 pro-$p$ 子群 $K\subset G_b(E)$，$(i^{b*}A)^K\in D(\Lambda)$ 自反。
\end{corollary}

\begin{proof}
该定理蕴含 $i^{b*}$ 与 $\mathbb{D}(\mathbb{D}(-))$ 交换，论证是归纳的。
\end{proof}

\subsection{ULA 层}

由于 $\Bun_G$ 是 Artin $v$-叠，对 $\pi\colon\Bun_G\to\ast$ 有 ULA 层的概念。（ULA 性在源上上同调光滑局部。）下面的命题是 ULA 性的「可对偶性」刻画的推论。

\begin{proposition}
$A\in D_{\et}(\Bun_G,\Lambda)$ 是 ULA 的，当且仅当
\[
p_1^*R\Hom(A,\Lambda)\otimes^L_\Lambda p_2^*A\to R\Hom(p_1^*A,p_2^*A)
\]
是同构，其中 $p_1,p_2\colon\Bun_G\times\Bun_G\to\Bun_G$ 是投影映射。
\end{proposition}

（先前陈述过一个用上 $!$ 的刻画；但由于 $\Bun_G$ 在一点上上同调光滑，它可以改用上 $*$ 表述。）

\begin{theorem}
$A\in D_{\et}(\Bun_G,\Lambda)$ 是 ULA 的，当且仅当对一切 $b\in B(G)$ 以及每个开 pro-$p$ 子群 $K\subset G_b(E)$，$(i^{b*}A)^{hK}\in D(\Lambda)$ 是完美复形。
\end{theorem}

\begin{proof}
我们需要判定 $\Bun_G\times\Bun_G$ 上某个层的映射是否为同构。这要用到我们对 $\Bun_G\times\Bun_G$ 上层范畴的知识。
\end{proof}

\begin{lemma}
外 $\boxtimes$-积
\[
(-)\boxtimes(-)\colon D_{\et}(\Bun_G,\Lambda)\boxtimes_{D(\Lambda)}D_{\et}(\Bun_G,\Lambda)\to D_{\et}(\Bun_G\times\Bun_G,\Lambda)
\]
\srcpage{112}%
是 $\infty$-范畴的等价。更确切地说，对 $A_1,A_2\in D_{\et}(\Bun_G,\Lambda)^\omega$，外张量积 $A_1\boxtimes A_2\in D_{\et}(\Bun_G\times\Bun_G,\Lambda)$ 紧，这类对象构成 $D_{\et}(\Bun_G\times\Bun_G,\Lambda)$ 的紧生成元；且对一切 $B_1,B_2\in D_{\et}(\Bun_G,\Lambda)$，
\[
R\Hom(A_1\boxtimes A_2,B_1\boxtimes B_2)\xleftarrow{\ \sim\ }R\Hom(A_1,B_1)\otimes^L_\Lambda R\Hom(A_2,B_2).
\]
\end{lemma}

\begin{proof}
分析紧生成元 $A^b_K$。
\end{proof}

\begin{proof}[续]
现在补完前述 ULA 判据（原文 Theorem 20.14）的证明。我们需要判定何时
\[
p_1^*R\Hom(A,\Lambda)\otimes p_2^*A\xrightarrow{\ \sim\ }R\Hom(p_1^*A,p_2^*A)
\]
是同构。对两边作用 $R\Hom(A_1\boxtimes A_2,-)$ 并用上述引理，就归结为证明
\[
R\Hom(\pi_\natural(A_1\otimes^L_\Lambda A),\Lambda)\otimes^L_\Lambda R\Hom(A_2,A)\to R\Hom(\pi_\natural(A_1\otimes^L_\Lambda A),R\Hom(A_2,A))
\]
是同构。它成立当且仅当 $\pi_\natural(A_1\otimes^L_\Lambda A)\in D(\Lambda)$ 完美。取 $A_1=i_{b!}\mathrm{c}\text{-}\mathrm{Ind}^{G_b(E)}_K\Lambda$ 并利用前述那一串等式（原文编号 (20.2.1)），即可看出这转化为 $(i^{b*}A)^K\in D(\Lambda)$ 完美。
\end{proof}
